\chapter*{LỜI NÓI ĐẦU}
\addcontentsline{toc}{chapter}{LỜI NÓI ĐẦU}

Trong bối cảnh nhu cầu sử dụng điện sinh hoạt ngày càng tăng, việc dự báo điện năng tiêu thụ của hộ gia đình có ý nghĩa quan trọng trong quản lý chi phí và điều chỉnh thói quen sử dụng điện.. Tuy nhiên, mức tiêu thụ điện của mỗi hộ không chỉ phụ thuộc vào số điện tháng trước mà còn chịu ảnh hưởng bởi số người sinh sống, thời điểm trong năm, vị trí địa lý và điều kiện khí tượng.

Đề tài ``Xây dựng hệ Neuro-Fuzzy thích nghi (ANFIS) dự báo điện năng tiêu thụ hộ gia đình dựa trên thói quen người dùng, thời gian, vị trí và yếu tố khí tượng'' được thực hiện nhằm nghiên cứu khả năng kết hợp giữa logic mờ và mạng nơ-ron trong bài toán dự báo. Hệ ANFIS có ưu điểm vừa học được từ dữ liệu, vừa có khả năng biểu diễn tri thức dưới dạng các luật IF--THEN dễ diễn giải.

Trong báo cáo này, nhóm trình bày cơ sở lý thuyết, quy trình chuẩn bị dữ liệu, thiết kế mô hình ANFIS, phương pháp đánh giá và định hướng triển khai hệ thống dự báo điện năng cho người dùng cuối. Do thời gian và dữ liệu còn hạn chế, báo cáo được xây dựng như một khung triển khai hoàn chỉnh để tiếp tục bổ sung kết quả thực nghiệm trong các giai đoạn sau.

Nhóm thực hiện xin chân thành cảm ơn giảng viên phụ trách học phần đã hướng dẫn và góp ý trong quá trình thực hiện đề tài.

\begin{flushright}
  \textit{Nhóm thực hiện}
\end{flushright}
